\section{相关工作}
\label{sec:related}

\paragraph{手动设计的核。}许多现有框架，如 TensorFlow XLA~\cite{tensorflow_xla, Tensorflow}、PyTorch~\cite{pytorch} 和 TensorRT~\cite{tensorrt}，依赖 GPU 专家为机器学习运算符手动设计核。
近期，大量工程努力被投入到手动优化常用 DNN（尤其是基础模型~\cite{bommasani2022opportunities}）的 GPU 核中。
例如，为加速注意力计算~\cite{transformers}，基于 FlashAttention~\cite{dao2023flash, tri2023flashdecoding, hong2024flashdecoding, fastertransformer} 开发了多种专用核。
由于现代 GPU 日益复杂——例如 A100 中的张量核心~\cite{Markidis2018tensorcores} 和 H100 中的线程块簇~\cite{nvidia-h100}——手动设计的核可能会错过难以手动发现的微妙优化。

\paragraph{基于超级优化的方法。}超级优化最初被引入以寻找最优指令序列~\cite{massalin, stoke, peephole}。
近期工作将超级优化技术应用于张量程序~\cite{TASO, wang2021pet, zheng2023einnet, tensat, unger2022unity, FlexFlow, hu2024korch, jeon2025graphpipe}。
然而，所有这些工作都只考虑核层次的代数变换，无法发现需要在核、块和线程所有层次联合考虑代数变换和调度变换的更复杂优化。
我们的评估表明，\sys 大幅优于现有 DNN 超级优化器，证明了多层次联合优化的重要性。

\paragraph{基于调度的方法。}近期工作引入了能自动优化核 GPU 执行调度的机器学习编译器。TVM~\cite{tvm, tvm_auto_tuner}、Ansor~\cite{ansor} 和 Triton~\cite{tillet2019triton} 等系统及其他工作~\cite{zheng2020flextensor, hagedorn2023graphene, feng2022tensorir}，均基于 Halide 中引入的算法-调度分离思想，搜索在 GPU 上执行用户指定算法的优化调度。
然而，基于调度的方法要求用户明确指定每个核的算法，其性能受限于所提供算法的质量。

\paragraph{多层次图表示。}Welder~\cite{shi2023welder} 和 ASPEN~\cite{park2023aspen} 引入了多层次块图，与 \sys 的 \graphs 有相似之处，因为两种表示都遵循 GPU 层次结构。然而，先前工作侧重于调度变换，而 \sys 通过同时考虑代数变换和发现新的自定义核，超越了单纯的调度优化。本文中提出的大多数优化方案超出了这些先前方法的范畴。
